% Generated draft: requires mathematical and source review.
\documentclass[11pt]{article}
\usepackage[margin=25mm]{geometry}
\usepackage{fontspec}
\setmainfont{Latin Modern Roman}
\usepackage{amsmath,amssymb,mathtools,unicode-math}
\setmathfont{Latin Modern Math}
\usepackage{enumitem,needspace,xurl,hyperref}
\setlength{\parindent}{0pt}
\setlength{\parskip}{5pt}
\setlength{\emergencystretch}{3em}
\newcommand{\N}{\mathbb{N}}
\newcommand{\Z}{\mathbb{Z}}
\newcommand{\Q}{\mathbb{Q}}
\newcommand{\R}{\mathbb{R}}
\newcommand{\C}{\mathbb{C}}
\newcommand{\F}{\mathbb{F}}
\newcommand{\A}{\mathbb{A}}
\newcommand{\PP}{\mathbb{P}}
\newcommand{\E}{\mathbb{E}}
\newcommand{\eps}{\varepsilon}
\newcommand{\dd}{\,\mathrm{d}}
\begin{document}
\section*{UnsolvedMath ID: 11000006}
\textbf{Source problem code:} catalogue entry
\subsection{Definitions and mathematical statement}
Problem 3.2 asks to find a \((6g-8)\)-obstructor complex \(L\) and a proper expanding map
\[
 F:L\times[0,\infty)\longrightarrow\mathcal T_g^{\ge\eps}.
\]
The source's obstructor-complex definition, expanding-map convention, allowed genus and \(\eps\), and target metric are not supplied in this extract. The printed product domain is retained; it is not silently replaced by a cone quotient.

In the usual surface setting, \(\mathcal T_g\) is Teichmüller space of marked hyperbolic structures on a closed oriented genus-\(g\) surface, and \(\mathcal T_g^{\ge\eps}\) denotes the thick subset whose shortest essential closed geodesic has length at least \(\eps>0\). This conventional meaning and its genus hypotheses require source confirmation. A map is proper when inverse images of compact sets are compact. One common simplicial expansion condition, for a specified metric \(d\), requires for disjoint simplices \(\sigma,\tau\subset L\)
\[
 \lim_{t\to\infty}
 \operatorname{dist}_d\bigl(F(\sigma\times[t,\infty)),
                           F(\tau\times[t,\infty))\bigr)=\infty.
\]
This is only an explanation of a standard convention, not a claim that the source's exact technical condition has been recovered. Likewise an \(m\)-obstructor is the technical object of the cited action-dimension framework; an arbitrary \(m\)-dimensional complex or an arbitrary nonembeddable complex is not an adequate replacement.

The supplied summary describes this as a construction strategy toward \(n(\operatorname{Mod}_g)=6g-6\), citing Bestvina--Kapovich--Kleiner and Bestvina--Feighn context, but it does not define \(n\) or supply a completed argument. The extraction also contains footnote 3 saying each sphere is triangulated as a join of \(0\)-spheres and \(L\) is triangulated as a join. Its sphere factors and antecedent are missing. The trailing page number and ``M. Bestvina'' text are page furniture, not additional hypotheses.
\subsection{Short English statement}
Construct the specified obstructor and proper expanding map into thick Teichmüller space. The technical definitions, metric and parameter ranges are missing. The supplied status assessment treats this as an uncompleted construction strategy, while the extracted footnote and page footer do not identify the desired complex.
\subsection{Sources}
\begin{itemize}
\item Local UnsolvedMath registry, permanent ID \(11000006\), exact source code \texttt{catalogue entry}, Problem 3.2, PDF page 14. No individual proposer or proposed year is supplied; the title is truncated.
\item Farb, edited volume \emph{Problems on Mapping Class Groups and Related Topics}: \url{https://www.math.uchicago.edu/~farb/papers/mcgbook.pdf}. The source has 329 explicitly typeset Problem/Conjecture/Question headings across contributed chapters, with local numbering disambiguated by PDF page. The extract includes footnote 3 and a footer naming M. Bestvina.
\item Supplied research summary describes a Farb (2006) strategy following Bestvina--Kapovich--Kleiner and Bestvina--Feighn toward an action-dimension bound; no complete bibliography entries or technical definitions are included.
\item Extraction pdf-font-structure, accessed 2026-07-29. Source evidence NEEDS\_REVIEW presents an open problem and requires a current primary-literature audit. Registry status open and classification OPEN-TRIAGE; supplied summary says no published completion of this exact construction was located.
\item Dataset attribution: Ulam AI / UnsolvedMath Contributors, CC BY 4.0. Source rights NEEDS\_REVIEW: publicly accessible author/publisher copy, redistribution terms requiring review.
\end{itemize}
\end{document}
